(1) 取向下挠度为正并忽略剪切变形，控制方程 $EI\,y^{(4)}=q$；固支边界 $y(0)=y'(0)=0$，自由端零弯矩与零剪力边界 $y''(L)=y'''(L)=0$。这是四阶边值问题，必须有四个条件，不能只使用端点位移条件。
(2) 四次积分和代入边界值给 $y(x)=q x^2(6L^2-4Lx+x^2)/(24EI)$。于是端点挠度 $y(L)=qL^4/(8EI)$，端点斜率 $y'(L)=qL^3/(6EI)$，二者均随刚度增加而减小，且挠度具长度量纲。
(3) 曲率相关弯矩大小为 $|M(x)|=EI y''(x)=q(L-x)^2/2$，剪力大小为 $|V(x)|=EI|y'''(x)|=q(L-x)$；固定端承受竖向反力 $qL$、弯矩 $qL^2/2$。载荷积分直接给总力 $qL$ 和关于根部的力矩 $q\int_0^Lx\,dx=qL^2/2$，与微分解一致。
(4) 弯曲能量 $U=\int_0^L M(x)^2/(2EI)\,dx=q^2L^5/(40EI)$。若从零逐渐按比例施加静载，外界做功为 $\frac12\int_0^Lq\,y(x)\,dx$；由 $\int_0^Ly(x)\,dx=qL^5/(20EI)$ 可得同一结果 $q^2L^5/(40EI)$，表明边界力学与能量公式相容。
